\chapter{Discussion and Conclusion}
This Master's thesis investigated the reduction in parameter count for \gls{DiT}s to enable more efficient image generation and facilitate their deployment in resource-constrained environments, particularly as new models become increasingly large and computationally intensive. An in-depth analysis of the compression framework was conducted on PixArt-$\Sigma$ (Sec. \ref{sec:PixArtExp}), evaluating several design choices and transferring the optimal settings to Flux.1-dev (Sec. \ref{sec:ExpFlux}). Surprisingly, the analysis revealed that low-rank approximation via truncated \gls{SVD} is a highly efficient method for compressing the projection matrices in individual transformer blocks while maintaining strong performance. This contrasts with earlier studies on \gls{LLM}s, which demonstrated that ordinary low-rank approximation led to significant quality loss and more advanced methods were required. Therefore, an importance-aware, rank-adaptive progressive compression approach, \textit{SVDtrunc}, was developed, establishing a new quality-parameter frontier for Flux.1-dev.


In previous structural pruning methods, model components were either removed entirely or replaced by a smaller module. In contrast, \gls{SVD}trunc leaves the overall architecture intact by performing a low-rank approximation of the weight matrices. Sec. \ref{sec:CompressionStrategyPixart} underlines that low-rank approximation provides a robust initialization for compressed models, leading to superior results compared to complete block pruning. This suggests that the redundancies in \gls{DiT}s are primarily located within individual transformer blocks rather than across them. It appears that each block fulfills its own functional task. Thus, by removing an entire block, the neighboring blocks must compensate for this loss, potentially limiting their capacity to perform their own functions. Even in a training-free approach (Sec. \ref{sec:trainingFreeCompression}) the \gls{SVD} compressed models outperform EcoDiff \cite{EcoDiff} and the block pruning strategy for models smaller than 90\% of their original size. This demonstrates that \gls{SVD}trunc better preserves the essential functionalities of each block. Particularly when retaining only 70\% of the original parameters, EcoDiff and block pruning collapse completely while \gls{SVD}trunc still generates viable images, even though a degradation in image quality is observed.

The analysis further revealed that the importance-aware block selection approach outperforms random or uniform blocks selection (Sec. \ref{sec:FluxBlockImportanceAnalysis}). Moreover, instead of compressing each selected block to a uniform size, a decreasing linear compression schedule was employed. By compressing less important blocks more aggressively, this approach demonstrated superior performance (Sec. \ref{sec:CompressionSchedulingPixart}). These results suggest that different blocks process tasks of varying complexity. Consequently, they exhibit different levels of resilience to compression, before a performance degradation occurs.

The effectiveness of \gls{SVD}trunc was demonstrated on Flux.1-dev (Sec. \ref{sec:FluxSchnellMainResults}).  \gls{SVD}trunc-m which is compressed to 68\% of the original model size of Flux.1-dev performs nearly as well as Flux.1-dev, only an average degradation of 0.75\% across all benchmarks is observed. As the model size was reduced further, the overall benchmark performance dropped gradually and predictably, such that even \gls{SVD}trunc-s outperforms competing methods with significantly higher model sizes. 

However, notwithstanding the superior benchmark performance and the comparable reduction in \gls{VRAM} consumption, a notable disadvantage of \gls{SVD}trunc is that the inference time decreases only marginally due to the increased number of sequential matrix operations. To truly establish a compression method for efficient \gls{DiT}s, the approach must be compatible with step-distilled models, as they require significantly fewer inference steps. Sec. \ref{sec:FluxSchnellMainResults} demonstrates that for compression ratios up to 30\%, the performance degradation of Flux.1-schnell is minor ($R=2.72\%$). However, when compressing Flux.1-schnell to 50\% of its original size, a significant performance drop occurs ($R=10.69\%$). The resulting images appear washed out and exhibit lower quality. This indicates that step-distilled models may contain fewer parameter redundancies and therefore cannot be compressed as aggressively as Flux.1-dev. The washed-out appearance closely resembles images generated by ordinary \gls{DiT}s using an insufficient number of steps. Consequently, in future work, it could be investigated whether it is beneficial to reverse the order by compressing the model first and subsequently applying step distillation. Alternatively, exploring a joint schedule integrating iterative compression and step-distillation training could be beneficial to better recover the model's performance.

The investigation of Flux.1-dev uncovered two distinct phenomena. First, a gradual domain shift was occasionally observed across different compression ratios, for instance from anime to realistic styles. Second, a loss of memorization regarding highly distinct concepts, such as "Bugs Bunny", became apparent while reducing model size. Notably, in both scenarios, even strongly compressed models (e.g., up to 39\% of the original size) continue to generate plausible, aesthetically pleasing images. This demonstrates that truncated \gls{SVD} is very effective in preserving parameters necessary for general image generation. However, fine details such as style specific nuances or highly specialized concepts appear to be stored in low-magnitude singular values that are removed and cannot be recovered. Consequently, \gls{SVD}trunc acts as a semantic filter that prioritizes structural coherence over detail. This filtration effect suggests that while \gls{SVD}trunc is ideal for general-purpose deployment where memory efficiency is paramount, specialized applications requiring high fidelity to specific artistic styles or intellectual properties might require more targeted preservation strategies

Future research could explore ``importance-aware`` truncation strategies, where singular values are preserved not based on their magnitude alone but on their contribution to specific high-value concepts or stylistic tokens. A potentially promising approach is \gls{GRASP} which precisely demonstrated in the training-free case that selecting singular values based on task importance rather than magnitude enables generation of images of greater detail. Therefore, it would be interesting to see if this advantage is also observed after subsequent fine-tuning.

Beyond importance-aware singular value selection, another avenue for future research could tackle the discrepancy between theoretical FLOPs and practical inference time. To reduce the actual latency, the development of hardware-optimized CUDA kernels specifically designed for low-rank matrix multiplication could be investigated. Furthermore, it would be interesting to explore whether \gls{SVD}trunc can be combined with low-bit quantization to further push the boundaries of deployment on consumer-grade hardware with even stricter \gls{VRAM} limitations.